\section{Related Work}
\label{sec:related_work}

Traffic Signal Control (TSC) has transitioned from civil engineering heuristics to decentralized artificial intelligence. This section surveys classical baselines, modern deep reinforcement learning (DRL) frameworks, recent benchmarking platforms, and the emerging literature on Emergency Vehicle Preemption (EVP).

\subsection{Classical and Optimal Signal Control Baselines}
Conventional traffic infrastructure relies on deterministic formulations. Webster~\cite{webster1958} established the foundational methodology for calculating \textbf{fixed-cycle phase splits} based on historical arrival distributions. To introduce responsiveness, adaptive schemes such as \textbf{SCATS}~\cite{lowrie1990scats} and \textbf{SCOOT}~\cite{hunt1981scoot} dynamically extend green allocations using inductive loop sensors.

In network traffic theory, the \textbf{Max-Pressure controller} developed by Varaiya~\cite{varaiya2013maxpressure} serves as the theoretical gold standard. Max-Pressure selects phases that maximize the pressure differential between incoming queues and downstream exit queues, mathematically guaranteeing network stabilization under stable demand regimes without prior knowledge of vehicle arrival rates. Modern benchmarking studies (e.g., LibSignal~\cite{mei2022libsignal}) consistently mandate comparison against Max-Pressure to establish whether learning-based policies exceed non-parametric mathematical heuristics.

\subsection{Deep Reinforcement Learning in Traffic Signal Control}
With the introduction of \textbf{Deep Q-Networks (DQN)}~\cite{mnih2015human}, reinforcement learning has been extensively applied to Intelligent Transportation Systems. Tabular multi-agent methods such as \textbf{MARLIN-ATSC}~\cite{eltantawy2013marlin} established decentralized coordination, while deep architectures like \textbf{IntelliLight}~\cite{wei2018intellilight} and \textbf{PressLight}~\cite{wei2019presslight} incorporated state representations that explicitly encode pressure and queue lengths.

More recently (2022--2025), attention-based multi-agent architectures—such as \textbf{CoLight}~\cite{wei2019colight} and \textbf{MPLight}~\cite{chen2020mplight}—have utilized Graph Neural Networks (GNNs) and spatio-temporal attention to model dependencies across adjacent intersections. However, as highlighted by benchmark studies like LibSignal~\cite{mei2022libsignal} and Sim2Signal~\cite{mei2024sim2signal}, two fundamental deficiencies persist in existing DRL traffic frameworks:
\begin{enumerate}
    \item \textbf{Absence of Physical Safety Bounds:} Pure end-to-end RL models operate with unconstrained action intervals. During early exploration or distribution shift, agents frequently output dangerously brief green windows ($<5$\,s) that induce rear-end collisions or violate pedestrian safety minimums.
    \item \textbf{Inconsistent Evaluation Baselines:} Numerous contemporary studies benchmark exclusively against naive, un-actuated fixed timers, failing to evaluate whether learning algorithms genuinely surpass industrial actuated controllers or theoretical Max-Pressure heuristics.
\end{enumerate}

\subsection{Emergency Vehicle Preemption (EVP) and Post-Clearance Shockwaves}
Emergency Vehicle Preemption (EVP) protocols have traditionally operated through optical strobes or Vehicle-to-Infrastructure (V2I) / Dedicated Short-Range Communications (DSRC)~\cite{klein2001sensor}. Recent advances (2023--2025) investigate edge-computing frameworks and multi-modal sensor fusion, utilizing acoustic siren recognition via Mel-Frequency Cepstral Coefficients (MFCCs)~\cite{salamon2014urbansound} paired with edge vision detectors (e.g., YOLOv4/v11~\cite{bochkovskiy2020yolov4}).

A critical challenge identified in recent 2024--2025 literature~\cite{zhang2024fogevp, liu2023collaborative} is the \textbf{preemption-congestion trade-off}: while forcing green waves significantly accelerates emergency transit (typically reducing transit latency by 40\% to 60\%), it abruptly starves conflicting arterial corridors, inducing acute secondary congestion shockwaves. Existing preemption protocols largely treat the emergency event as an isolated disruption, failing to model the \textbf{Post-Preemption Recovery Time ($T_{\text{rec}}$)} required to restore intersection equilibrium. AURA directly resolves this gap by coupling deterministic safety envelopes with a momentum clearing reward that actively suppresses post-clearance queue spikes.
